\subsection{Detailed direct certificates for the nonzero-gap terminals}
\label{subsec:new-detailed-direct-certificates}

This subsection supplies the calculations suppressed in the preceding case
proofs.  They are written as inequalities for actual reaches and radial
witnesses.  There is no formal chain notation.

\subsubsection{The disk--gap support calculation}

Let $c=c_{\max}(p,q)$, $m=\min\{p,q\}$, and
$d=1-c$.  The regular radial hexagon has inradius $hd$.  Choose the farther
endpoint $G$ of the complementary gap.  Its norm is
\[
 \rho=\sqrt{1-m+m^2}.
\]
For a unit normal whose angle from $G$ is $\theta$, the three point supports
are
\[
 \rho\cos\theta,
 \quad
 \rho\cos(\theta+2\pi/3),
 \quad
 \rho\cos(\theta+4\pi/3).
\]
The disk replaces every one of these by at least $hd$.  A minimizing
orientation is either a disk-only orientation or a tie orientation.  At a tie
one projection equals $hd$.  The other active projection is
\[
 \frac{-hd+\sqrt{3(\rho^2-h^2d^2)}}2.
\]
Thus the support sum is at least $h$ exactly when
\[
 3d(1-d)\ge m(1-m).
 \tag{6.51}
\]
Since $d=1-c$, this is (6.29).

We now derive the local inequality used in the high-$c$ part.  Work in a
support orientation with angular parameter $t\in[0,\pi/3]$ and put
\[
 x=\cos\left|t-\frac\pi6\right|.
\]
The support inequalities for the local kite
$\conv\{0,pu,qv,c(u+v)\}$ imply
\[
 cx\le h
 \tag{6.52}
\]
and, after taking the smaller boundary demand $m$,
\[
 mx+\frac c2
 \left(x+\sqrt3\sqrt{1-x^2}\right)\le h.
 \tag{6.53}
\]
When $c>h$, the interval is $h\le x\le h/c$.  The left side of (6.53) is
concave.  If $c>1-m$, its value at $x=h$ is $h(m+c)>h$, so feasibility forces
its value at the other endpoint to be at most $h$:
\[
 \frac mc+\frac12+\sqrt{c^2-\frac34}\le1.
\]
Hence
\[
 m\le c\left(\frac12-\sqrt{c^2-\frac34}\right).
 \tag{6.54}
\]
Substitution of the right side into (6.51) gives
\[
\begin{aligned}
&3c(1-c)-
 c\left(\frac12-\sqrt{c^2-\frac34}\right)
 \left[1-c\left(\frac12-\sqrt{c^2-\frac34}\right)\right]\\
&\quad=
\frac{c(1-c)}2
\left(5-2c-2c^2+2\sqrt{c^2-\frac34}\right)\ge0.
\end{aligned}
\]
This proves the complementary-gap inequality directly from the four local
support conditions.

\subsubsection{The CE2 one-gap threshold inequalities}

We prove the two inequalities in (6.39).  Put
\[
 X=R-\delta,
 \qquad
 Q=\frac{\eta+\alpha+\delta}{2R}.
\]
The right-trace inequality gives
\[
 WX=RW-W\delta=\eta+P-W\delta>\eta+\alpha.
 \tag{6.55}
\]
The selected low root $\epsilon(\alpha)$ is the smaller root of
\[
 z^2-(1-\alpha)z+(1-\alpha)^2-(1-\alpha)^4=0.
 \tag{6.56}
\]
It is enough to prove that $X$ lies above that root.  Since
$0<X<R$ and the quadratic opens upward, substitute the lower estimate in
(6.55).  Equivalently consider
\[
 \pi(q)=(\eta+q)(1-2q)-2qW.
\]
This is concave.  At the endpoints of the admissible interval $0\le q\le P$,
\[
 \pi(0)=\eta>0,
\]
while the identities $RW=\eta+P$ and $P=E\eta$ give
\[
 \pi(P)=\eta(\eta+2ER^2)>0.
\]
Thus $\pi(\alpha)>0$, which is precisely
\[
 \epsilon(\alpha)<X.
 \tag{6.57}
\]

For the second inequality put $T=\alpha+\delta$.  Multiply
\[
 \alpha+W\delta<P,
 \qquad
 R\alpha+\delta<P
\]
by $W$ and $R$, respectively, and add.  Since
$W+R^2=W^2+R=E^2$,
\[
 E^2T<P=E\eta,
 \qquad
 T<\frac\eta E.
 \tag{6.58}
\]
Let $d=\min\{\alpha,\delta\}$.  Then $d\le T/2$.  By (6.5),
\[
 \min\{\epsilon(\alpha),\epsilon(\delta)\}
 <\frac{2d}{1-2d}
 \le\frac{T}{1-T}.
\]
It remains to prove $T/(1-T)<Q$.  This is equivalent to
\[
 2RT<(\eta+T)(1-T).
\]
Define
\[
 \chi(T)=(\eta+T)(1-T)-2RT.
\]
The function is concave.  At $T=0$, $\chi(0)=\eta>0$.  At
$T=\eta/E$, direct simplification gives
\[
 \chi\left(\frac\eta E\right)
 =\frac{\eta}{E^2}(E-R)^2>0.
\]
By (6.58), $T$ lies strictly between these endpoints.  Hence
\[
 \boxed{
 \min\{\epsilon(\alpha),\epsilon(\delta)\}<Q.
 }
 \tag{6.59}
\]
This proves the CE2 dichotomy without a composed reach map.

\subsubsection{The CE2 short-ray root separation}

We give the remaining details of Theorem~\ref{thm:new-ce2-short-ray}.  The
identities
\[
 RW=\eta+P,
 \qquad P=E\eta
\]
and the CE2 inequalities imply
\[
 Wq=W(R-\delta)>\eta+\alpha,
\]
\[
 Rp=R(W-\alpha)>\eta+\delta.
\]
Thus
\[
 p,q>\eta+e.
 \tag{6.60}
\]
Let $M=\max\{R,W\}$.  Since $\alpha,\delta\ge e$,
\[
 e<\frac{P}{1+M}.
\]
Moreover
\[
 (1+M)^2-3E^2=(2M-1)(2-M)\ge0,
\]
so
\[
 \eta>\sqrt3e.
 \tag{6.61}
\]
The universal bound
\[
 P\le\frac{2\sqrt3-3}{4},
 \qquad 1+M\ge\frac32
\]
gives
\[
 e<\frac{2\sqrt3-3}{6}=e_*.
\]

For $c=1-e$, the roots of
$f_c(t)=c^4-c^2+ct-t^2$ are $t_\pm$ from the proof above.  At
$t_0=(1+\sqrt3)e$,
\[
 f_c(t_0)=eB(e),
\]
with
\[
 B(e)=e^3-4e^2-3\sqrt3e+\sqrt3-1.
\]
Since
\[
 B'(e)=3e^2-8e-3\sqrt3<0
\]
on $[0,e_*]$, one has
\[
 B(e)\ge B(e_*)=\frac{302\sqrt3-501}{72}>0.
\]
Therefore $t_-<t_0<p,q$.  Also
\[
 p+q=1-\alpha-\delta\le1-2e=c-e.
\]
If, say, $q\le p$, then $q>t_-$ and
\[
 p<c-e-t_-<c-t_-=t_+.
\]
The same holds with $p,q$ interchanged.  Hence
\[
 t_-<p,q<t_+,
\]
which proves $f_c(p),f_c(q)>0$.  In the exact local finite-caliper formula at
radial demand $c$, the four side candidates are
\[
 p+c,
 \quad q+c,
 \quad\frac{c^2}{\sqrt{p^2-pc+c^2}},
 \quad\frac{c^2}{\sqrt{q^2-qc+c^2}}.
\]
The first two exceed one because $p,q>e$, and the last two exceed one exactly
because $f_c(p),f_c(q)>0$.  This completes the direct short-ray proof.

\subsubsection{Detailed adjacent Vd residual estimate}

We justify (6.48).  In the adjacent placement define the residuals
$\rho_R,\rho_L$ as in Case 1 of Proposition~\ref{prop:new-one-vd-assembly}.
The signed CE2 domain first gives
\[
 \alpha+\delta<\frac16\min\{R,W\}.
 \tag{6.62}
\]
The positive residual $\rho_L$ has one of two forms.

If
\[
 \rho_L=W-\alpha,
\]
then
\[
 4\delta+\alpha
 \le4(\alpha+\delta)
 <\frac{2W}{3}<W,
\]
so $4\delta<\rho_L$.

Assume instead
\[
 \rho_L=1-A_0.
\]
The other residual cannot equal $1-B_0$, since (6.47) would imply
$A_0+B_0>3/2$, contrary to the endpoint-distance bound
$A_0+B_0\le2/\sqrt3$.  Hence
\[
 \rho_R=1-(W+\delta),
 \qquad
 B_0\ge y:=\frac{k}{R}.
\]
The two points of $T_0$ with reaches $1-\rho_L$ and $y$ have distance at most
one.  The diameter-transfer inequality gives
\[
 \rho_L>\frac y2>\frac{\eta}{2R}.
 \tag{6.63}
\]
Since $\rho_R+\rho_L<1/2$,
\[
 \delta>R-\frac12+\frac{\eta}{2R}=:J(R).
 \tag{6.64}
\]
Using $\eta/R=W/(1+E)$, exact differentiation gives
\[
4E(1+E)^2J'(R)
=2E(1+E)(1+2E)-W(2R-1)>0.
\]
Thus $J$ is strictly increasing.  Since $J(3/8)=1/24$ and (6.62) gives
$\delta<1/24$, equation (6.64) forces $R<3/8$.  Then
\[
 (2-3R)^2-E^2=(3-8R)(1-R)>0,
\]
so $E<2-3R$ and
\[
 \frac{\eta}{R}=\frac{W}{1+E}>\frac13.
\]
Together with (6.63),
\[
 \rho_L>\frac16>4\delta.
\]
This proves (6.48) in both residual cases.

\subsubsection{Detailed nonadjacent Vd separation}

Let $x=k/W$ and $y=k/R$ be the two center near endpoints.  Let $B,U$ be the
farthest already-covered extents on the corresponding edges.  We prove
\[
 B\le M_0(x),
 \qquad
 U\le M_0(y).
 \tag{6.65}
\]
By symmetry it is enough to treat $B$.  The residual component formula says
that $B$ is either $W+\delta$ or the supercritical endpoint $B_0$.
If $B=W+\delta$, the center contains boundary points with reaches $x,B$, so
the diameter-one inequality gives $B\le M_0(x)$.  If $B=B_0$ and
$A_0<x$, then the opposite covered extent is $U=A_0$; equation (6.47) would
again force $A_0+B_0>3/2$.  Hence $A_0\ge x$, and monotonicity of $M_0$
gives
\[
 B_0\le M_0(A_0)\le M_0(x).
\]
This proves (6.65).

The elementary identity
\[
 1-M_0(z)>\frac z2
 \qquad(z>0)
\]
then gives
\[
 \rho_R>\frac x2,
 \qquad
 \rho_L>\frac y2.
\]
The common center endpoint-slack bound is
\[
 \alpha+\delta<\frac12\min\{x,y\}.
\]
Thus (6.49) follows.  The exact Vd corner form gives the strict own-radial
bound used in Case 2.

\subsubsection{Detailed Vd1 rescue inequalities}

Use the Vd1 corner parameter $t\ge1$, let
$d=\sqrt{t^2+t+1}$, and write the supported interval as $[c,u]$.  The endpoint
relations are
\[
 c=\frac{t(1-b)}{t+1},
 \qquad
 u=\frac{d-a-tb-1}{t}.
\]
The midpoint condition gives
\[
 a+tb\le d-1-\frac t2.
\]
Eliminating $tb=t-c(t+1)$ gives
\[
 a\le F(t,c):=d-1-\frac{3t}{2}+c(t+1).
\]
For fixed $c\le1/2$,
\[
 \frac{\partial F}{\partial t}
 =\frac{2t+1}{2\sqrt{t^2+t+1}}-\frac32+c<0,
\]
so
\[
 a\le F(1,c)=\sqrt3-\frac52+2c=:L(c).
\]
The inequality $2L(c)\le1-M_c^{\rm sup}$ is equivalent, after squaring, to
\[
 Q(c)=20c^2+(18\sqrt3-52)c+47-24\sqrt3\ge0.
\]
The polynomial decreases on $[0,1/2]$ and
$Q(1/2)=26-15\sqrt3>0$.  Hence
\[
 a\le1-M_c^{\rm sup}.
\]

Put $\varepsilon=1-u$.  Direct substitution gives
\[
 \varepsilon=\varepsilon_0+\frac at,
 \qquad
 \varepsilon_0=
 \frac{2t+1-c(t+1)-d}{t}>0.
\]
The function
\[
 z\longmapsto\frac{z}{z+\varepsilon_0+z/t}
\]
is increasing.  Since
\[
 \varepsilon_0+\frac{F(t,c)}t=\frac12,
\]
we obtain
\[
 \frac{a}{a+1-u}
 \le\frac{F(t,c)}{F(t,c)+1/2}
 \le2L(c)
 \le1-M_c^{\rm sup}.
\]
This proves the two direct Vd1 rescue inequalities in (6.50).
